\documentclass[conference]{IEEEtran}
\IEEEoverridecommandlockouts

\usepackage{cite}
\usepackage{amsmath,amssymb,amsfonts}
\usepackage{graphicx}
\usepackage{textcomp}
\usepackage{xcolor}
\usepackage{booktabs}
\usepackage{array}
\usepackage[utf8]{inputenc}
\usepackage[T1]{fontenc}
\graphicspath{{figs/svm_linear/}}

\def\BibTeX{{\rm B\kern-.05em{\sc i\kern-.025em b}\kern-.08em
    T\kern-.1667em\lower.7ex\hbox{E}\kern-.125emX}}

\begin{document}

\title{Classificação de Áreas Queimadas em Imagens Landsat\\
Utilizando SVM Linear (Support Vector Machine)}

\author{
  \IEEEauthorblockN{Pedro (Sobrenome)}
  \IEEEauthorblockA{
    \textit{Reconhecimento de Padrões}\\
    Trabalho Final\\
    Recife, Brasil
  }
}

\maketitle

\begin{abstract}
Este trabalho apresenta uma metodologia de reconhecimento de padrões para
detecção de áreas queimadas utilizando imagens Landsat e índices de
vegetação NDVI obtidos em períodos pré e pós-fogo. O método emprega um
classificador SVM linear treinado a partir de amostras balanceadas de pixels
queimados e não queimados provenientes de múltiplas cenas. Foram utilizadas
variáveis derivadas do NDVI, incluindo diferença temporal, razão espectral e
variação absoluta, visando aumentar a capacidade discriminativa do modelo.
O treinamento foi realizado com duas cenas Landsat e posteriormente aplicado
em uma terceira cena independente para validação. O modelo apresentou boa
capacidade de generalização espacial, sendo avaliado por métricas como
acurácia, precisão, recall, F1-score, IoU e matriz de confusão. Além disso,
foi gerado um mapa temático classificado das áreas queimadas detectadas.
\end{abstract}

\begin{IEEEkeywords}
área queimada, SVM linear, NDVI, sensoriamento remoto, Landsat,
classificação binária, desbalanceamento de classes.
\end{IEEEkeywords}

%=====================================================================
\section{Introdução}

A detecção automática de áreas queimadas em imagens de satélite é uma tarefa
de alta relevância para o monitoramento ambiental, gestão territorial e
estudos do ciclo do carbono. Métodos baseados em limiares espectrais como o
dNBR são amplamente utilizados~\cite{b2}, porém dependem de calibração manual
e não exploram a capacidade discriminativa de múltiplas variáveis combinadas.

O objetivo deste trabalho é desenvolver um classificador binário baseado em
SVM linear (\textit{LinearSVC}) capaz de detectar automaticamente pixels
queimados e não queimados a partir de atributos derivados do índice NDVI em
imagens Landsat~8/9. O modelo é treinado em duas cenas e validado em uma
terceira cena independente, avaliando a generalização espacial do classificador.

%=====================================================================
\section{Dados de Entrada e Saída}

\subsection{Entradas}

O sistema utiliza como dados de entrada:

\begin{itemize}
  \item \textbf{Imagens raster Landsat} contendo:
    \begin{itemize}
      \item NDVI pré-queimada (\texttt{ndvi\_pre})
      \item NDVI pós-queimada (\texttt{ndvi\_pos})
    \end{itemize}
  \item \textbf{Máscaras binárias de treinamento}:
    \begin{itemize}
      \item \texttt{0} = área não queimada
      \item \texttt{1} = área queimada
    \end{itemize}
  \item Conjunto de cenas Landsat para treinamento e teste.
\end{itemize}

\subsection{Saídas}

O processamento gera:

\begin{itemize}
  \item Modelo SVM treinado;
  \item Predição das classes queimado/não queimado;
  \item Métricas quantitativas de desempenho: Accuracy, Precision, Recall,
        F1-score e IoU;
  \item Matriz de confusão;
  \item Mapa classificado final em formato GeoTIFF.
\end{itemize}

%=====================================================================
\section{Metodologia}

\subsection{Pré-processamento e Organização dos Dados}

O experimento foi desenvolvido a partir de imagens Landsat previamente
processadas para geração do índice de vegetação NDVI em períodos pré e
pós-queimada. Para cada cena analisada foram utilizados:

\begin{itemize}
  \item NDVI anterior ao evento de fogo;
  \item NDVI posterior ao evento;
  \item máscara binária de referência contendo pixels classificados como
        queimados (1) e não queimados (0).
\end{itemize}

Os rasters foram carregados utilizando a biblioteca \texttt{Rasterio},
sendo realizada a leitura apenas da banda principal de cada arquivo GeoTIFF.
Como diferentes cenas podem apresentar pequenas discrepâncias dimensionais,
os rasters foram ajustados para dimensões compatíveis por meio do recorte ao
menor tamanho comum entre as imagens pré e pós-evento.

Posteriormente, foi realizado um procedimento de limpeza numérica para
eliminar inconsistências frequentes em produtos de sensoriamento remoto:
valores nulos (\texttt{NaN}), infinitos e pixels inválidos. Além disso,
regiões associadas a corpos d'água foram removidas do conjunto de
treinamento por meio de máscara baseada em valores negativos de NDVI,
reduzindo falsos positivos causados pela semelhança espectral com áreas
queimadas. Por fim, os atributos foram limitados ao intervalo espectral
$[-1, 1]$, reduzindo efeitos de valores extremos e garantindo estabilidade
numérica durante o treinamento.

\subsection{Balanceamento e Amostragem dos Dados}

O forte desbalanceamento entre classes — pixels não queimados muito mais
frequentes que queimados — pode induzir o modelo supervisionado a favorecer
excessivamente a classe majoritária. Para minimizar esse problema, foi
adotada uma estratégia de amostragem balanceada por classe: os pixels de
cada classe foram separados individualmente e foi realizada uma seleção
aleatória estratificada com a mesma quantidade de amostras para ambas as
classes.

O número máximo de amostras por classe foi definido em 25.000 pixels,
totalizando aproximadamente 50.000 amostras no treinamento. Essa estratégia
apresentou três vantagens principais: redução do viés para a classe dominante,
maior estabilidade estatística do treinamento e redução do custo computacional
do SVM. O balanceamento contribuiu significativamente para melhorar métricas
associadas à classe minoritária, especialmente Recall e F1-score.

\subsection{Treinamento e Validação do Classificador SVM}

O modelo supervisionado utilizado foi o Support Vector Machine (SVM),
implementado por meio da classe \texttt{LinearSVC} da biblioteca
\texttt{scikit-learn}. Antes do treinamento, os atributos espectrais foram
normalizados utilizando o \texttt{StandardScaler} (média zero, variância
unitária) — etapa fundamental em modelos SVM, pois o algoritmo é sensível
à escala dos atributos.

O treinamento foi realizado utilizando duas cenas Landsat independentes,
enquanto uma terceira cena foi reservada exclusivamente para validação do
modelo, avaliando a capacidade de generalização espacial em áreas não
utilizadas durante o treinamento. O modelo foi configurado com:

\begin{itemize}
  \item penalização linear (\texttt{LinearSVC});
  \item ponderação balanceada entre classes
        (\texttt{class\_weight='balanced'});
  \item limite elevado de iterações para convergência numérica.
\end{itemize}

A função objetivo do SVM consiste em encontrar um hiperplano ótimo capaz de
maximizar a separação entre as classes no espaço multidimensional de atributos
espectrais. Após o treinamento, o modelo foi aplicado pixel a pixel sobre a
cena de teste para geração das predições temáticas.

%=====================================================================
\section{Resultados}

\subsection{Métricas de Desempenho}

O modelo apresentou acurácia global de 80,14\%. Contudo, devido ao forte
desbalanceamento entre as classes, essa métrica isoladamente não representa
adequadamente o desempenho do classificador. A Tabela~\ref{tab:metrics}
resume os resultados obtidos na cena de teste.

\begin{table}[!t]
  \caption{Métricas de desempenho — Cena de Teste}
  \label{tab:metrics}
  \centering
  \begin{tabular}{lc}
    \toprule
    \textbf{Métrica} & \textbf{Valor} \\
    \midrule
    Acurácia (OA)    & 80,14\% \\
    Precision        & 14,04\% \\
    Recall           & 83,43\% \\
    F1-score         & 24,04\% \\
    IoU              & 13,66\% \\
    \bottomrule
  \end{tabular}
\end{table}

O elevado Recall (83,43\%) indica alta capacidade de detecção da classe
queimada, reduzindo a ocorrência de falsos negativos. Em contrapartida, a
baixa Precision (14,04\%) evidencia elevada incidência de falsos positivos,
refletida também nos valores de IoU (13,66\%) e F1-score (24,04\%).

\subsection{Matriz de Confusão}

A Tabela~\ref{tab:cm} apresenta a matriz de confusão obtida na cena de teste.

\begin{table}[!t]
  \caption{Matriz de Confusão — Cena de Teste (Cena~3)}
  \label{tab:cm}
  \centering
  \begin{tabular}{lcc}
    \toprule
    & \textbf{Pred. Não Queimado} & \textbf{Pred. Queimado} \\
    \midrule
    \textbf{Real Não Queimado} & 38.501 (TN) & 9.616 (FP) \\
    \textbf{Real Queimado}     &    312 (FN) & 1.571 (TP) \\
    \bottomrule
  \end{tabular}
\end{table}

O \textit{Classification Report} completo é apresentado na
Tabela~\ref{tab:classreport}.

\begin{table}[!t]
  \caption{Classification Report — Cena de Teste}
  \label{tab:classreport}
  \centering
  \begin{tabular}{lcccc}
    \toprule
    \textbf{Classe} & \textbf{Precision} & \textbf{Recall}
                    & \textbf{F1} & \textbf{Support} \\
    \midrule
    Não queimado (0) & 0,99 & 0,80 & 0,89 & 48.117 \\
    Queimado (1)     & 0,14 & 0,83 & 0,24 & 1.883  \\
    \midrule
    Accuracy         & \multicolumn{4}{c}{0,80 (50.000 amostras)} \\
    Macro avg        & 0,57 & 0,82 & 0,56 & 50.000 \\
    Weighted avg     & 0,96 & 0,80 & 0,86 & 50.000 \\
    \bottomrule
  \end{tabular}
\end{table}

A Fig.~\ref{fig:cm} ilustra visualmente a distribuição dos erros.
A matriz de confusão demonstra que o modelo assumiu um comportamento mais
sensível à detecção de queimadas, priorizando a identificação da classe
minoritária mesmo ao custo de superestimativas espaciais.

\begin{figure}[!t]
  \centering
  \includegraphics[width=0.85\columnwidth]{fig_cm.png}
  \caption{Matriz de Confusão — Cena~3 (teste). TN=38.501, FP=9.616,
           FN=312, TP=1.571.}
  \label{fig:cm}
\end{figure}

%=====================================================================
\section{Discussão}

Os resultados confirmam que o modelo SVM linear é capaz de generalizar a
detecção de áreas queimadas entre diferentes cenas Landsat, apresentando
elevada sensibilidade para identificação da classe de interesse (Recall~$= 83\%$).
Apesar da boa capacidade de detecção, observou-se significativa ocorrência de
falsos positivos, principalmente em regiões com comportamento espectral
semelhante às áreas queimadas — como solos expostos, pastagens degradadas
e sombras de nuvem.

A estratégia de balanceamento por amostragem igualitária (25.000 amostras por
classe) foi fundamental para viabilizar a detecção da classe minoritária,
evitando que o classificador ignorasse sistematicamente os pixels queimados.
Contudo, a limitação a apenas dois atributos de entrada (\texttt{ndvi\_pre} e
\texttt{ndvi\_pos}) restringe a capacidade discriminativa do modelo, visto que
solos expostos e áreas queimadas podem apresentar valores de NDVI similares.

A incorporação de atributos adicionais como as bandas espectrais brutas (B4,
B5, B7) e o NBR multitemporal poderiam melhorar a separabilidade entre classes
e reduzir os falsos positivos.

%=====================================================================
\section{Conclusão}

O modelo SVM linear demonstrou capacidade de generalização espacial para
detecção de áreas queimadas em imagens Landsat, com Recall de 83,43\% na
cena de teste. A metodologia mostrou potencial para aplicações em
monitoramento ambiental, podendo ser aprimorada com:

\begin{enumerate}
  \item incorporação de novos atributos espectrais (bandas brutas, NBR);
  \item técnicas de pós-processamento espacial para redução de falsos positivos;
  \item expansão do conjunto de treinamento com mais cenas e períodos distintos.
\end{enumerate}

A abordagem de amostragem balanceada mostrou-se eficaz para lidar com o
desbalanceamento de classes inerente ao problema de detecção de queimadas.

\begin{thebibliography}{00}

\bibitem{b1}
V.~N.~Vapnik,
``The Nature of Statistical Learning Theory,''
Springer, New York, 1995.

\bibitem{b2}
C.~H.~Key e N.~C.~Benson,
``Landscape Assessment: Composite Burn Index,''
in \textit{FIREMON}, USDA Forest Service, RMRS-GTR-164-CD, 2006.

\bibitem{b3}
E.~Chuvieco \textit{et al.},
``Generation and analysis of a new global burned area product based on
MODIS 250~m reflectance bands,''
\textit{Earth System Science Data}, vol.~11, pp.~529--551, 2019.

\bibitem{b4}
F.~Pedregosa \textit{et al.},
``Scikit-learn: Machine Learning in Python,''
\textit{J. Mach. Learn. Res.}, vol.~12, pp.~2825--2830, 2011.

\end{thebibliography}

\end{document}
